\section{Literature review}
\label{sec:literature}

A large body of work on earlier general-purpose technologies documents that technological change over recent decades has been skill-biased and, more precisely, routine-biased: computerisation substituted for workers performing codifiable, routine tasks in the middle of the wage distribution while complementing non-routine cognitive work at the top \citep{autor2003}. The distributional consequence was a widening of wage inequality and a polarisation of employment. Crucially, however, studies that trace these labour market shocks through to household disposable income find that tax-benefit systems absorbed much of the shock: progressive taxation and means-tested transfers substantially muted the pass-through of rising wage dispersion into household income inequality \citep{doorley2023automation}. This distinction between market-income and disposable-income effects motivates the microsimulation approach adopted in this paper.

Artificial intelligence appears to break with the routine-biased pattern in an important respect. Occupation-level exposure measures constructed from the capabilities of modern AI systems---whether based on linkages between AI applications and occupational abilities \citep{felten2021}, overlap between patent text and job task descriptions \citep{webb2019}, or assessments of which tasks large language models can perform \citep{eloundou2023}---consistently show that exposure rises with education and earnings. Professional, managerial and technical occupations score among the most exposed, a finding replicated for European labour markets \citep{williamson2024}. Yet high exposure need not imply displacement: the same workers who are most exposed may also enjoy the largest complementarity gains, since AI can augment rather than replace complex cognitive work. \citet{cazzaniga2024} formalise this ambiguity, showing that the net effect of AI on inequality depends on the balance between displacement and complementarity within highly exposed, high-income occupations, and can plausibly go in either direction.

Direct empirical evidence on the labour market effects of AI adoption remains limited and, where available, mixed. Firm-level studies generally find that AI adopters are larger, more productive and faster-growing than non-adopters, and that they pay higher wages, though it is difficult to separate the causal effect of adoption from selection into adoption by already-successful firms. Worker-level evidence is similarly inconclusive: some studies find employment growth at adopting firms alongside compositional shifts towards more skilled workers, while others detect reductions in hiring for exposed roles. This thin and heterogeneous evidence base is one reason scenario-based simulation, of the kind pursued here and in \citet{doorley2026}, remains a useful complement to ex-post estimation.

A related strand quantifies AI's productivity potential at the task level. \citet{eloundou2023} estimate that large language models, with appropriate complementary software, could enable roughly 15 per cent of all US worker tasks to be completed significantly faster at the same quality, with the share rising markedly once tooling built on top of the models is taken into account. Experimental and quasi-experimental studies of specific occupations point in a consistent direction: in programming, professional writing and customer support, access to generative AI tools raises output per worker, with the largest gains typically accruing to less experienced or lower-skill workers, thereby compressing within-occupation productivity gaps. Whether this within-job compression translates into narrower wage dispersion depends on how the gains are shared and on which jobs survive.

Evidence specific to the United Kingdom is beginning to accumulate. \citet{kleinteeselink2025} link measures of generative AI exposure to UK firm-level employment records and find that more exposed firms reduced employment by around 4.5 per cent following the diffusion of generative AI, with losses concentrated in junior positions, which fell by roughly 5.8 per cent. From a macro-distributional perspective, \citet{rockall2025} apply the IMF's exposure-and-complementarity framework to the UK and conclude that, in contrast to earlier automation waves, AI could reduce wage inequality if displacement effects fall disproportionately on high-income workers---though sufficiently strong complementarity among those same workers could offset or reverse this compression. Our paper speaks directly to this open question by tracing both scenarios through the UK tax-benefit system.

The concentration of adjustment among junior workers found for the UK echoes emerging evidence of seniority-biased effects elsewhere. Using US resume and job-posting data, \citet{hosseini2026} document that employment of junior workers falls by around 9 per cent within six quarters of firms' AI adoption, while senior employment is essentially unaffected, consistent with AI substituting for entry-level task content while complementing experienced judgement. In a similar vein, \citet{brynjolfsson2025} find a 13 per cent relative decline in employment among workers aged 22--25 in the most AI-exposed occupations since late 2022, with older workers in the same occupations largely spared. If AI's incidence is graded by seniority as much as by skill, the distributional consequences differ from those of a classic skill-biased shock, since junior workers are spread across the earnings distribution.

Finally, AI may reshape the distribution of income between factors as well as within labour. Task-based models of automation imply that the substitution of capital for labour in a widening range of tasks lowers the labour share and channels a growing portion of national income to the owners of capital \citep{acemoglu2018}. Because capital income is far more concentrated than labour income, this channel raises household inequality even if wage dispersion is unchanged, and it strengthens the case examined by \citet{bastani2024} for progressive taxation of capital income as a policy response. Our simulations therefore consider not only wage scenarios but also the sensitivity of UK disposable-income inequality to shifts in the capital share.
